\clearpage
\cleardoublepage

\chapter{Modern Training Techniques}

\section{Introduction}

Training modern deep learning systems involves far more than gradient descent. This chapter covers advanced training techniques that enable training on massive datasets, transfer knowledge across tasks, and achieve state-of-the-art performance. These techniques are essential for anyone working with contemporary neural networks.

\section{Pre-training and Fine-tuning}

The dominant paradigm for building powerful models is pre-training followed by fine-tuning.

\subsection{Pre-training Objectives}

Pre-training learns general representations from large unlabeled data:

\begin{itemize}
    \item \textbf{Language Modeling (CLM):} Predict the next token given previous tokens
    \item \textbf{Masked Language Modeling (MLM):} Predict randomly masked tokens (BERT-style)
    \item \textbf{Denoising Autoencoder:} Reconstruct corrupted input
    \item \textbf{Next Sentence Prediction (NSP):} Predict sentence ordering
\end{itemize}

For vision:
\begin{itemize}
    \item \textbf{Contrastive Learning:} Maximize similarity between positive pairs
    \item \textbf{Masked Autoencoding (MAE):} Reconstruct masked patches
    \item \textbf{Self-distillation:} Student learns from teacher
\end{itemize}

\subsection{Fine-tuning Strategies}

After pre-training, models are adapted to specific tasks:

\begin{enumerate}
    \item \textbf{Full fine-tuning:} Update all parameters on target task data
    \item \textbf{Linear probing:} Train only a linear classifier on frozen features
    \item \textbf{Adapter methods:} Add small trainable modules between layers
    \item \textbf{LoRA (Low-Rank Adaptation):} Optimize low-rank decompositions of weight matrices
\end{enumerate}

\begin{note}[When to use each approach]
\begin{itemize}
    \item \textbf{Abundant task data:} Full fine-tuning typically best
    \item \textbf{Limited task data:} LoRA/adapters prevent overfitting
    \item \textbf{Many tasks:} Adapter methods enable parameter-efficient multi-task learning
\end{itemize}
\end{note}

\section{LoRA: Low-Rank Adaptation}

LoRA reduces trainable parameters by decomposing weight updates:

\begin{equation}
\mathbf{W} \leftarrow \mathbf{W} + \Delta\mathbf{W}, \quad \Delta\mathbf{W} = \mathbf{B}\mathbf{A}
\end{equation}

where $\mathbf{B} \in \mathbb{R}^{d \times r}$, $\mathbf{A} \in \mathbb{R}^{r \times k}$, and $r \ll \min(d, k)$.

During forward pass with input $\mathbf{x}$:
\begin{equation}
\mathbf{h} = \mathbf{W}\mathbf{x} + \frac{\alpha}{r}\mathbf{B}\mathbf{A}\mathbf{x}
\end{equation}

where $\alpha$ is a scaling factor (typically equal to rank $r$).

\begin{code}
    \captionof{listing}{LoRA Implementation in PyTorch}
    \label{code:lora}
\begin{lstlisting}[language=python]
import torch
import torch.nn as nn

class LoRALinear(nn.Module):
    def __init__(self, in_features, out_features, rank=4, alpha=4):
        super().__init__()
        self.rank = rank
        self.alpha = alpha
        self.scaling = alpha / rank
        
        # Original frozen weights
        self.weight = nn.Parameter(
            torch.randn(out_features, in_features), 
            requires_grad=False
        )
        
        # Trainable low-rank matrices
        self.lora_A = nn.Parameter(torch.randn(rank, in_features))
        self.lora_B = nn.Parameter(torch.zeros(out_features, rank))
        
        # Initialize A with random Gaussian, B with zeros
        nn.init.normal_(self.lora_A, std=1.0 / rank)
    
    def forward(self, x):
        # Original forward pass
        h = F.linear(x, self.weight)
        # LoRA update
        h += self.scaling * (x @ self.lora_A.T @ self.lora_B.T)
        return h

# Usage in a model
class LoRALinearLayer(nn.Module):
    def __init__(self, linear_layer, rank=4, alpha=4):
        super().__init__()
        self.linear = linear_layer
        self.lora = LoRALinear(
            linear_layer.in_features,
            linear_layer.out_features,
            rank=rank,
            alpha=alpha
        )
    
    def forward(self, x):
        return self.lora(x)
\end{lstlisting}
\end{code}

Trainable parameters reduced from $d \times k$ to $r \times (d + k)$.

\section{Transfer Learning}

Transfer learning leverages knowledge from related tasks.

\subsection{Feature Transfer}

Features learned on one domain often transfer well:

\begin{itemize}
    \item ImageNet features transfer to medical imaging, satellite imagery
    \item BERT features transfer across NLP tasks
    \item Key insight: Lower layers capture general features; higher layers capture task-specific patterns
\end{itemize}

\subsection{Progressive Freezing}

Gradually unfreeze layers during fine-tuning:
\begin{enumerate}
    \item Train only the classifier head (frozen backbone)
    \item Unfreeze last $n$ layers, fine-tune
    \item Optionally unfreeze remaining layers
\end{enumerate}

This prevents catastrophic forgetting of pre-trained features.

\section{Multi-Task Learning}

Training a single model on multiple tasks simultaneously.

\subsection{Hard Parameter Sharing}

Share encoder weights across tasks:
\begin{equation}
\mathcal{L}_{\text{total}} = \sum_{t=1}^{T} \lambda_t \mathcal{L}_t(\theta_{\text{shared}}, \theta_t)
\end{equation}

where $\lambda_t$ is task weight and $\theta_t$ are task-specific heads.

\subsection{Soft Parameter Sharing}

Each task has its own parameters, with regularization encouraging similarity:
\begin{equation}
\mathcal{L}_{\text{total}} = \sum_{t=1}^{T} \mathcal{L}_t(\theta_t) + \lambda \sum_{i,j} \|\theta_i^{(l)} - \theta_j^{(l)}\|^2
\end{equation}

\subsection{Task Balancing}

Multi-task learning requires careful task weighting:
\begin{itemize}
    \item \textbf{Equal weighting:} Simple but may not optimal
    \item \textbf{Uncertainty weighting:} Weight by inverse task uncertainty (Kendall et al., 2018)
    \item \textbf{Gradient balancing:} Balance gradient magnitudes across tasks
    \item \textbf{Dynamic weighting:} Adjust weights during training
\end{itemize}

\section{Continual Learning}

Training sequentially on new tasks without forgetting previous ones.

\subsection{Catastrophic Forgetting}

Neural networks tend to forget previous tasks when trained on new ones. Mitigation strategies:

\begin{enumerate}
    \item \textbf{Experience Replay:} Store samples from previous tasks and replay during new training
    \item \textbf{Elastic Weight Consolidation (EWC):} Protect important weights with regularization
    \item \textbf{Progressive Networks:} Add new columns for new tasks while freezing old ones
    \item \textbf{Architectural Growth:} Expand network capacity for new tasks
\end{enumerate}

\subsection{EWC Regularization}

\begin{equation}
\mathcal{L}_{\text{EWC}} = \mathcal{L}_B(\theta) + \frac{\lambda}{2}\sum_i \mathcal{F}_i (\theta_i - \theta_{A,i}^*)^2
\end{equation}

where $\mathcal{F}$ is the Fisher information matrix capturing parameter importance.

\section{Self-Supervised Learning}

Learning representations without labels using pretext tasks.

\subsection{Contrastive Learning}

Maximize agreement between differently augmented views of the same image:

\begin{equation}
\mathcal{L} = -\log \frac{\exp(\text{sim}(\mathbf{z}_i, \mathbf{z}_j)/\tau)}{\sum_{k=1}^{2N} \mathbf{1}_{[k \neq i]} \exp(\text{sim}(\mathbf{z}_i, \mathbf{z}_k)/\tau)}
\end{equation}

where $\tau$ is temperature, $\mathbf{z}$ are normalized embeddings.

Popular contrastive learning frameworks:
\begin{itemize}
    \item \textbf{SimCLR:} Simple framework with strong augmentations
    \item \textbf{MoCo:} Momentum contrast for large negative samples
    \item \textbf{SwAV:} Cluster assignments instead of pairwise comparison
\end{itemize}

\subsection{Masked Image Modeling}

BERT-style masking for images (MAE, BEiT):
\begin{itemize}
    \item Mask random patches (75\% for MAE)
    \item Decoder reconstructs masked patches from visible ones
    \item Encoder operates only on visible patches (highly efficient)
\end{itemize}

\section{Distillation and Compression}

Transferring knowledge from large models to smaller ones.

\subsection{Knowledge Distillation}

A smaller student model learns from a larger teacher:
\begin{equation}
\mathcal{L}_{\text{distill}} = \alpha \cdot \text{CE}(y, \text{softmax}(z_s/T)) + (1-\alpha) \cdot \text{CE}(y, \text{softmax}(z_t/T))
\end{equation}

where $z_s$ and $z_t$ are student and teacher logits, $T$ is temperature.

Key insight: Soft targets from teacher contain more information than hard labels.

\subsection{Self-Distillation}

A model distills knowledge into itself:
\begin{itemize}
    \item Train model A, then use A's predictions to train model B of same size
    \item BYOL, SimSiam: Two views with predictor network
    \item Shakes out noise and produces sharper predictions
\end{itemize}

\section{Large Batch Training}

Scaling training to thousands of GPUs requires careful batch size considerations.

\subsection{Linear Scaling Rule}

When increasing batch size $B$, scale learning rate proportionally:
\begin{equation}
\alpha_{\text{new}} = \alpha_{\text{old}} \times \frac{B_{\text{new}}}{B_{\text{old}}}
\end{equation}

Goyal et al. (2017) achieved 256-way data parallelism on ImageNet with 8K batch size.

\subsection{Warmup and Scaling}

For very large batches ($>$16K):
\begin{itemize}
    \item Gradual warmup to prevent early instability
    \item Layer-wise learning rate decay (LLRD): Lower LR for earlier layers
    \item Longer training (adjust schedule accordingly)
\end{itemize}

\section{Efficient Training Techniques}

\subsection{Gradient Checkpointing}

Trade compute for memory by recomputing activations during backward pass:
\begin{itemize}
    \item Store activations only at selected checkpoints
    \item Recompute forward pass segments during backward
    \item Memory reduction: $O(\sqrt{n})$ with $O(n)$ compute overhead
\end{itemize}

\subsection{Offloading}

Move tensors between CPU and GPU memory:
\begin{itemize}
    \item CPU offload: Keep parameters on CPU, load to GPU per layer
    \item ZeRO (Zero Redundancy Optimizer): Partition optimizer states across GPUs
    \item Activations offload: Pipeline recomputation
\end{itemize}

\section{Debugging Training}

\subsection{Common Issues and Solutions}

\begin{table}[ht]
    \centering
    \caption{Common Training Issues and Solutions}
    \begin{tabular}{L{3cm} L{4cm} L{4cm}}
        \toprule
        \textbf{Symptom} & \textbf{Possible Cause} & \textbf{Solution} \\
        \midrule
        Loss NaN & Large gradients & Reduce LR, gradient clipping \\
        Loss not decreasing & Wrong learning rate & LR range test, reduce LR \\
        Overfitting & Too few data, too many params & Regularization, augmentation \\
        Underfitting & Model too simple & Larger model, more features \\
        Catastrophic forgetting & Single-task training & EWC, replay, progressive \\
        Instability at init & Poor initialization & Kaiming/He init, proper normalization \\
        \bottomrule
    \end{tabular}
\end{table}

\subsection{Training Recipes}

Proven recipes for common scenarios:

\begin{algorithm}[H]
\caption{Training Recipe for Vision Models}
\begin{enumerate}
    \item Start with small learning rate (1e-4 for Adam, 0.1 for SGD)
    \item Use LR range test to find optimal LR
    \item Apply cosine annealing schedule with warmup
    \item Use strong augmentation (RandAugment, Mixup, CutMix)
    \item Label smoothing (0.1 for classification)
    \item Weight decay (0.01 for AdamW)
    \item Exponential moving average of weights (EMA, decay=0.9999)
    \item Early stopping on validation loss
\end{enumerate}
\end{algorithm}

\begin{algorithm}[H]
\caption{Training Recipe for LLMs}
\begin{enumerate}
    \item Pre-train on large corpus (trillion+ tokens for GPT-3 scale)
    \item Use mixed precision (BF16) with gradient checkpointing
    \item Apply gradient clipping (1.0)
    \item Use learning rate warmup (2000-10000 steps)
    \item Cosine decay to small LR (typically 10\% of peak)
    \item Weight decay (0.1 for LLaMA-2)
    \item Attention dropout during fine-tuning
    \item LoRA for efficient fine-tuning on downstream tasks
\end{enumerate}
\end{algorithm}

\section{Chapter Summary}

\begin{enumerate}
    \item Pre-training on large datasets followed by fine-tuning is the dominant paradigm
    \item LoRA enables parameter-efficient fine-tuning by decomposing weight updates
    \item Transfer learning allows leveraging representations across domains
    \item Multi-task learning requires careful task balancing strategies
    \item Continual learning techniques prevent catastrophic forgetting
    \item Self-supervised learning (contrastive, masked) learns from unlabeled data
    \item Knowledge distillation transfers capabilities from large to small models
    \item Large batch training benefits from linear learning rate scaling and warmup
    \item Systematic debugging and proven recipes accelerate development
\end{enumerate}

\section{Further Reading}

\begin{itemize}
    \item \cite{hinton2015distilling} --- Distilling the Knowledge in a Neural Network
    \item \cite{devlin2018bert} --- BERT: Pre-training of Deep Bidirectional Transformers
    \item \cite{kendall2018multi} --- Multi-Task Learning Using Uncertainty to Weigh Losses
    \item \cite{kirkpatrick2017overcoming} --- Overcoming Catastrophic Forgetting in Neural Networks
    \item \cite{hu2021lora} --- LoRA: Low-Rank Adaptation of Large Language Models
    \item \cite{chen2020simple} --- A Simple Framework for Contrastive Learning (SimCLR)
    \item \cite{he2022masked} --- Masked Autoencoders Are Scalable Vision Learners (MAE)
    \item \cite{goyal2017accurate} --- Accurate, Large Minibatch SGD: Training ImageNet in 1 Hour
    \item \cite{rombach2022high} --- High-Resolution Image Synthesis with Latent Diffusion Models
    \item \cite{wei2022regard} --- Chain-of-Thought Prompting Elicits Reasoning in Large Language Models
\end{itemize}
